% Partie : histoire
% Chapitre : biographies
% Section : Alan Turing

\documentclass[../../../include/open-logic-section]{subfiles}

\begin{document}

\olfileid[fr]{his}{bio}{tur}

\olsection{Alan Turing}

Alan Turing naquit à Maida Vale, à Londres, le 23 juin 1912. Il est
considéré comme le père de l'informatique théorique. Son intérêt pour
les sciences physiques et les mathématiques apparut dès son plus
jeune âge. Mais ses écoles accordaient surtout de l'importance à la
littérature et aux humanités classiques, et faisaient peu de place à
ses centres d'intérêt. Ses résultats scolaires furent donc médiocres,
et nombre de ses enseignants le réprimandèrent.

\olphoto{turing-alan}{Alan Turing}

Turing étudia les mathématiques au King's College de Cambridge.
En 1936, il conçut ce qu'on appelle aujourd'hui la machine de Turing
pour définir précisément la notion de fonction calculable et
démontrer l'indécidabilité du problème de la décision. Alonzo Church
l'avait devancé en établissant ce résultat au moyen de son propre
lambda-calcul. L'article de Turing fut néanmoins publié, avec une
référence au résultat de Church. Celui-ci invita Turing à Princeton,
où il séjourna de 1936 à 1938 et obtint un doctorat sous sa direction.

Malgré son intérêt pour la logique, Turing conserva son goût pour
les sciences physiques. Ses compétences pratiques furent mises à
profit au service britannique de cryptanalyse de Bletchley Park
pendant la Seconde Guerre mondiale. Il joua un rôle central dans le
déchiffrement du code Enigma utilisé par les communications navales
allemandes. Ses connaissances en statistique et en cryptographie,
associées à l'introduction de machines électromécaniques, permirent à
l'équipe de déchiffrer le code grâce à une machine appelée
«~bombe~». Les travaux de Turing sur le chiffre de Lorenz
contribuèrent aussi aux méthodes employées à Bletchley Park.
C'est pour ce chiffre allemand que Tommy Flowers conçut Colossus,
le premier ordinateur électronique programmable au monde.

Turing était homosexuel. En 1941, il demanda pourtant en mariage
Joan Clarke, l'une de ses collègues à Bletchley Park, et lui confia
son homosexualité avant de rompre leurs fiançailles. Il eut
plusieurs partenaires au cours de sa vie, alors que les actes
homosexuels étaient encore des infractions pénales au Royaume-Uni.
En 1952, un ami de son compagnon de l'époque cambriola sa maison.
Lorsqu'il déposa plainte, Turing reconnut entretenir une relation
homosexuelle, pensant que le gouvernement s'apprêtait à légaliser
ces relations. Ce n'était pas le cas, et il fut poursuivi pour
«~gross indecency~» (indécence grave). Pour éviter la prison, il
choisit un traitement hormonal destiné à réduire sa libido. Turing
fut retrouvé mort le 8 juin 1954 à la suite d'une surdose de cyanure,
probablement par suicide. La reine Élisabeth~II lui accorda une
grâce royale en 2013.

\begin{reading}
Pour une biographie approfondie d'Alan Turing, voir
\citet{Hodges2014}. Sa vie et ses travaux ont inspiré la pièce
\emph{Breaking the Code}, adaptée pour la télévision en 1996 avec
Derek Jacobi dans le rôle de Turing. Le film \emph{The Imitation
Game}, nommé aux Oscars et interprété notamment par Benedict
Cumberbatch et Keira Knightley,
s'inspire aussi librement de la vie de Turing et de son séjour à
Bletchley Park \citep{Imitation2014}.

\citet{Radiolab2012} propose plusieurs balados sur la vie et les
travaux de Turing. Le documentaire de BBC Horizon
\emph{The Strange Life and Death of Dr.~Turing} est visible en ligne
\citep{Sykes1992}. \citep{Theelen2012} présente une courte vidéo
d'une machine de Turing en LEGO qui fonctionne, réalisée pour le
centenaire de sa naissance en 2012.

L'article original de Turing sur les machines de Turing et le
problème de la décision est \citet{Turing1937}.
\end{reading}

\end{document}
